% 地球（综述）
% license CCBYSA3
% type Wiki

本文根据 CC-BY-SA 协议转载翻译自维基百科\href{https://en.wikipedia.org/wiki/Earth}{相关文章}。

\begin{figure}[ht]
\centering
\includegraphics[width=6cm]{./figures/5280dfaa768e0550.png}
\caption{“蓝色弹珠”，阿波罗 17 号，1972 年 12 月} \label{fig_Earth_1}
\end{figure}
地球是距太阳的第三颗行星，也是唯一已知能够孕育生命的天体。这得益于地球是一颗海洋世界——在太阳系中，唯一能够维持液态地表水的星球。地球上几乎所有的水都存在于全球性海洋中，覆盖了地球地壳的 70.8\%。剩余的 29.2\% 地壳为陆地，其中大部分集中在地球陆半球的大陆地块上。地球的大多数陆地区域至少具有一定湿度并被植被覆盖，而地球极地沙漠中的大型冰盖所储存的水量超过了地球地下水、湖泊、河流和大气水的总和。地球的地壳由缓慢移动的构造板块组成，彼此相互作用形成山脉、火山与地震。地球拥有液态外核，可产生磁层，能够偏转大部分具有破坏性的太阳风和宇宙射线。

地球拥有一个动态的大气层，这维持了地球的表面环境，并在陨石体及紫外线入射时提供保护。大气层主要由氮气和氧气组成。水汽在大气层中广泛存在，形成覆盖地球大部分区域的云层。水汽是一种温室气体，并与大气中的其他温室气体（特别是二氧化碳 $CO_2$）一道，通过捕获来自太阳光的能量，创造出液态表面水和水汽能够持续存在的条件。这一过程维持了目前 14.76 °C（58.57 °F）的平均地表温度，在正常大气压下，水在该温度呈液态。不同地理区域所捕获能量的差异（如赤道地区获得的阳光多于极地地区）驱动了大气和海洋环流，形成具有不同气候区的全球气候系统，并产生降水等多种天气现象，使碳和氮等组成成分得以循环。

地球呈近似椭球形，周长约为 40,000 公里（24,900 英里）。它是太阳系中密度最高的行星。在四颗类地行星中，地球最大且质量最高。地球距太阳约八光分（1 天文单位），沿轨道绕太阳运行一周约需一年（约 365.25 天）。地球自转一周所需时间略少于一天（约 23 小时 56 分钟）。地球自转轴相对于其绕太阳轨道平面的垂线有倾斜，从而产生四季更替。地球拥有一颗永久天然卫星——月球。月球以 384,400 公里（238,855 英里，约 1.28 光秒）的距离绕地球运行，其直径约为地球的四分之一。月球的引力有助于稳定地球自转轴，产生潮汐，并逐渐减缓地球自转。同样，地球的引力已使月球自转潮汐锁定，月球永远以同一面朝向地球。

地球与太阳系中大多数其他天体一样，约在 45 亿年前由早期太阳系中的气体和尘埃形成。在地球历史的最初十亿年中，海洋形成，随后生命在海洋中出现。生命在全球扩散，并持续改变地球的大气和地表，导致约 20 亿年前发生“大氧化事件”。人类约 30 万年前在非洲出现，随后扩散到地球的每一块大陆。人类依赖地球的生物圈和自然资源而生存，但对地球环境的影响却日益增加。如今，人类对地球气候与生物圈的影响已不可持续，威胁着人类自身及许多其他生命形式的生计，并导致物种广泛灭绝。

地球呈椭球形，周长约 40,000 公里（24,900 英里）。它是太阳系中密度最高的行星。在四颗岩质行星中，地球最大且质量最高。地球距离太阳约八光分（1 个天文单位），绕太阳公转一周约需一年（约 365.25 天）。地球绕自身轴心自转一周时间略少于一天（约 23 小时 56 分钟）。地球自转轴相对于其绕太阳轨道平面垂线存在倾斜，从而产生四季变化。地球拥有一颗永久天然卫星——月球，月球以 384,400 公里（238,855 英里，即 1.28 光秒）的距离绕地运行，其直径约为地球的四分之一。月球的引力有助于稳定地球自转轴、引发潮汐，并逐渐减缓地球自转。同样，地球的引力已使月球自转产生潮汐锁定，月球始终以同一面朝向地球。

地球与太阳系中大多数其他天体一样，约在 45 亿年前由早期太阳系中的气体和尘埃形成。在地球历史最初的十亿年中，海洋开始形成，随后生命在其中诞生。生命在全球扩散，并持续改变地球的大气层和地表，导致约二十亿年前发生“大氧化事件”。人类约在 30 万年前出现在非洲，并扩散至地球的每一块大陆。人类依赖地球的生物圈和自然资源生存，但对地球环境的影响日益加剧。目前人类对地球气候和生物圈的影响已不可持续，威胁着人类和许多其他生命形式的生存，并正在引发广泛的物种灭绝。
\subsection{词源}
现代英语单词 earth 经由中古英语发展而来，源自古英语名词 eorðe $^{[22]}$。在所有日耳曼语言中都有其同源词，由此重建出原始日耳曼语形式 erþō。在最早的文献中，单词 eorðe 用来翻译拉丁语 terra 和希腊语 gē 的多重含义：地面、泥土、陆地、人类世界、世界表面（包括海洋）以及整个地球本身。与罗马神话中的 Terra（或 Tellus）和希腊神话中的 Gaia 类似，“地球”可能在日耳曼多神教中也被视为一位拟人化的女神：晚期北欧神话中出现过 Jörð（“大地”），这位女巨人常被认为是雷神托尔的母亲 $^{[23]}$。

历史上，“earth”通常以小写书写。在中古英语早期，为表达“地球”这一明确含义，开始使用短语 the earth。到了早期现代英语时期，名词大写逐渐普及，the earth 也写作 the Earth，尤其是在与其他天体并列提及时。近来，这一名称有时直接写作 Earth，以与其它行星名称类比，然而 earth 和 the earth 的用法仍然常见 $^{[22]}$。不同出版风格各有规定：牛津拼写认为小写形式最为常见，而大写形式则为可接受的变体。另一种惯例是在作为名称使用时采用大写（如 “Earth's atmosphere”），但在前有定冠词时使用小写（如 “the atmosphere of the earth”）。在口语表达中几乎总是写作小写，例如“what on earth are you doing?” $^{[24]}$

名称 Terra（/ˈtɛrə/）有时用于科学写作；它也常见于科幻作品中，用来将人类居住的行星与其他星球区分开来 $^{[25]}$。而在诗歌中，Tellus（/ˈtɛləs/）则被用来表示地球的拟人形象 $^{[26]}$。Terra 也是一些拉丁语系语言（如意大利语和葡萄牙语）中“地球”的名称；在其他罗曼语中，该词的拼写略有变化，如西班牙语 Tierra 与法语 Terre。希腊诗意名称 Gaia 的拉丁化形式 Gaea 在英语中较为少见，尽管由于“盖亚假说（Gaia hypothesis）”，另一种拼写 Gaia 已变得常用，并读作 /ˈɡaɪ.ə/，而不是更传统的英语发音 /ˈɡeɪ.ə/ $^{[27]}$。

关于地球还有许多形容词。earthly（地上的、尘世的）源自 Earth。由拉丁语 Terra 衍生出 terran（/ˈtɛrən/）、terrestrial（/təˈrɛstriəl/），以及（经由法语） terrene（/təˈriːn/）；由拉丁语 Tellus 衍生出 tellurian（/tɛˈlʊəriən/）和 telluric $^{[28][29][30][31][32]}$。
\subsection{自然历史}
\subsubsection{形成}
\begin{figure}[ht]
\centering
\includegraphics[width=8cm]{./figures/409cc725593ef0b7.png}
\caption{对形成地球及其他太阳系天体的早期太阳系原行星盘的描绘} \label{fig_Earth_2}
\end{figure}
在太阳系中发现的最古老物质被测定为距今 $4.5682^{+0.0002}_{-0.0004}
$ 十亿年前$^{[33]}$。在距今 $4.54\pm0.04$ 十亿年前，原始地球已经形成$^{[34]}$。太阳系中的各个天体与太阳一起形成并演化。理论上，太阳星云通过引力坍缩从分子云中分割出一块空间，开始旋转并扁平化成一个环星盘，然后行星在该盘内与太阳一起生长。星云包含气体、冰粒和尘埃（包括原始核素）。根据星云理论，微行星通过吸积而形成，据估计原始地球可能需要约 7000 万年至 1 亿年才能形成$^{[35]}$。

对月球年龄的估计从距今 45 亿年到显著更年轻不等$^{[36]}$。主流假说认为，月球由一次撞击后从地球上抛出的物质吸积而成——一个约为地球质量 10\%、大小与火星相似的天体（命名为 忒伊亚 Theia）与地球发生碰撞$^{[37]}$。它以掠擦方式撞击地球，其部分质量与地球合并$^{[38][39]}$。大约在距今 40 亿年至 38 亿年前，晚期重轰炸**期间大量小行星撞击导致月球大范围表面环境发生显著改变，并据此推断地球当时的表面环境也发生了重大变化$^{[40]}$。
\subsubsection{形成之后}
地球的大气和海洋由火山活动和气体逸出形成$^{[41]}$。来自这些源的水蒸气凝结成海洋，并由小行星、原行星和彗星带来的水与冰进一步补充$^{[42]}$。按照一种模型，自地球形成之初起，地球上就可能已经有足以填满海洋的水量$^{[43]}$。在该模型中，大气中的温室气体阻止了当时只有当前 70\%光度的新生太阳导致海洋结冰$^{[44]}$。到距今约 35 亿年前，地球的磁场已经建立，这有助于防止大气被太阳风剥离$^{[45]}$。
\begin{figure}[ht]
\centering
\includegraphics[width=8cm]{./figures/e931dfc771eed8c4.png}
\caption{“淡橙色小点——对早期地球的一种印象画，呈现其富含甲烷并带有橙色色调的早期大气层$^{[46]}$。”} \label{fig_Earth_3}
\end{figure}
随着地球熔融的外层冷却，形成了第一批固体地壳，被认为成分为镁铁质$^{[47]}$。第一批大陆地壳成分更为长英质，是由这种镁铁质地壳部分熔融形成的$^{[47]}$。在太古宙沉积岩中发现的冥古宙时代锆石颗粒表明，至少部分长英质地壳早在 44 亿年前就已经存在，比地球形成仅晚约 1.4 亿年$^{[48]}$。

关于这一初始、体积很小的大陆地壳如何演化至如今的丰富状态，主要有两种模型$^{[49]}$：
(1) 相对稳定地持续增长直至今天$^{[50]}$，这一观点得到全球大陆地壳放射性定年数据的支持；
(2) 在太古宙期间大陆地壳体积迅速增长，形成了现今大陆地壳的主体$^{[51][52]}$，这一观点得到来自锆石中铪同位素和沉积岩中钕同位素证据的支持。通过大陆地壳的大规模循环再造，特别是在地球早期阶段，这两种模型及其证据是可以调和的$^{[53]}$。

新的大陆地壳是通过板块构造形成的，而板块构造最终由地球内部持续散失热量驱动。在数亿年的时间尺度上，构造力促使大陆地壳区域结合成超级大陆，随后又破裂分离。约在 7.5 亿年前，最早已知的超级大陆之一罗迪尼亚（Rodinia）开始裂解。之后大陆重新聚合，在 6 亿至 5.4 亿年前形成泛诺西亚（Pannotia），最终又形成盘古大陆（Pangaea），该超级大陆在约 1.8 亿年前开始裂解$^{[54]}$。

最近一次冰期循环模式始于约 4,000 万年前$^{[55]}$，并在约 300 万年前的更新世显著增强$^{[56]}$。自那以后，高纬度和中纬度地区经历了反复的冰川与间冰期循环，大约每 21,000 年、41,000 年和 100,000 年重复一次$^{[57]}$。最近一次冰期，俗称“最后冰期（Last Glacial Period）”，使大陆大部分区域至中纬度地区都被冰层覆盖，并于约 11,700 年前结束$^{[58]}$。
